\section{Conclusion}
\label{sec:conclusion}

This paper presented a characterization study of decoder computation in 3D medical image segmentation architectures. By evaluating voxel-level prediction shifts between matched full and efficient decoders, we established a fundamental distinction between decoder activity (where computation alters predictions) and decoder net gain (where computation improves overall correctness). Our results show that while useful decoder benefit is highly concentrated at organ boundaries ($K_{80} = 5\%$), global net gain remains net-neutral as positive corrections are offset by negative degradations. Furthermore, we showed that inference-time uncertainty signals reliably locate decoder activity but cannot predict directional decoder utility.

These findings demonstrate that prediction difficulty and computational utility are fundamentally distinct quantities, revealing the inherent limits of static decoding and naive uncertainty gating. Future work should focus on utility-aware dynamic decoders that jointly optimize spatial routing and feature refinement.
